% ============================================================
% CAPITOLO 6 — Progetto 3: Server REST API con Node.js
% ============================================================
\chapter{Progetto 3: Server REST API con Node.js}

\section*{Cosa Costruirai}
Un \textbf{server REST API} completo con Node.js e Express che:
\begin{itemize}
  \item Espone endpoint CRUD per la gestione di ``note''
  \item Valida tutti gli input con Zod\index{Zod}
  \item Gestisce errori con middleware centralizzato
  \item Genera documentazione Swagger automatica
  \item Include test di integrazione
\end{itemize}

Questo server diventerà la \textbf{base del progetto full-stack} che crescerà nei Capitoli~7--10.

\textbf{Tempo stimato}: 45--60 minuti

\begin{notebox}
\textbf{Box Teoria --- Cos'è una REST API?} Una REST API è un'interfaccia che permette a programmi diversi di comunicare tra loro tramite richieste HTTP standard. Funziona come un cameriere in un ristorante: il client (browser, app mobile) invia un \textit{ordine} (richiesta) specificando \textit{cosa vuole} (URL) e \textit{che operazione} (metodo HTTP: GET per leggere, POST per creare, PUT per aggiornare, DELETE per eliminare). Il server elabora la richiesta e risponde con i dati in formato JSON --- un formato testuale leggibile sia dagli umani che dai computer. Esempio: \texttt{GET /api/notes} chiede la lista delle note; il server risponde con \texttt{[\{"id": 1, "title": "Prima nota"\}]}.
\end{notebox}

\begin{notebox}
\textbf{Box Tooling --- Stack scelto per questo esempio.}
\begin{itemize}
  \item \textbf{Runtime:} Node.js~20~LTS
  \item \textbf{Framework:} Express.js~4.x
  \item \textbf{Validazione:} Zod~3.x
  \item \textbf{Test:} Jest~29.x
\end{itemize}
\textbf{Alternative equivalenti:} Python/FastAPI, Go/Gin, Java/Spring Boot, C\#/ASP.NET. Il \textbf{pattern architetturale} (endpoint CRUD, validazione input, middleware, test) è identico indipendentemente dal linguaggio. Se preferisci un altro stack, applica lo stesso metodo 0-code: il \texttt{\_CONTEXT.md} guiderà l'IA a generare l'implementazione nel framework che hai scelto.
\end{notebox}

% ---------------------------------------------------------
\section{Prerequisiti}

\begin{versionbox}
Le istruzioni di questo capitolo sono verificate con Node.js~20~LTS, Express.js~4.x, Zod~3.x e Jest~29.x. I framework JavaScript evolvono costantemente: se l'IA ti segnala una libreria deprecata o suggerisce una versione più recente (es.\ Express~5, Vitest al posto di Jest), fidati della correzione e procedi. I concetti architetturali restano identici.
\end{versionbox}

\begin{praticabox}[title={Pratica --- Installare Node.js}]
\textbf{Windows:} Scarica da \url{https://nodejs.org} la versione LTS (20.x+).

\textbf{macOS:}
\begin{lstlisting}[language=bash]
brew install node
\end{lstlisting}

\textbf{Linux:}
\begin{lstlisting}[language=bash]
curl -fsSL https://deb.nodesource.com/setup_20.x \
  | sudo -E bash -
sudo apt-get install nodejs
\end{lstlisting}

\textbf{Verifica:}
\begin{lstlisting}[language=bash]
node --version    # v20.x.x
npm --version     # 10.x.x
\end{lstlisting}
\end{praticabox}

% ---------------------------------------------------------
\section{Il Contesto del Progetto (ADLC Fase 0-2)}\index{REST API}\index{Express.js}\index{Node.js}

\begin{notebox}
\textbf{Nota didattica}: In un progetto reale, prima di scrivere il \texttt{\_CONTEXT.md} applicheresti la Fase~3 del workflow professionale (Sezione~3.8): l'IA propone le opzioni tecnologiche, tu scegli, e documenti tutto in \texttt{docs/DESIGN.md} come ADR. In questo progetto guidato, le scelte sono già state fatte per te --- così puoi concentrarti sull'imparare il metodo ADLC senza il sovraccarico decisionale. Nel progetto autonomo dell'Appendice~F applicherai il workflow completo, Design incluso.
\end{notebox}

\begin{praticabox}[title={Pratica --- Setup e \texttt{\_CONTEXT.md}}]
Crea la cartella \texttt{notes-api}, aprila in VS~Code e crea il file \texttt{\_CONTEXT.md}:
\end{praticabox}

\begin{lstlisting}[language={},title={\texttt{\_CONTEXT.md} per Notes API}]
# Progetto: Notes API

## Scopo
REST API per la gestione di note personali. Backend che verra'
usato come fondamento per un'app full-stack (React + Flutter).

## Tecnologie
- Runtime: Node.js 20 LTS
- Framework: Express.js 4.x
- Linguaggio: JavaScript ES2022+ con ESModules
- Validazione: Zod
- Documentazione: swagger-jsdoc + swagger-ui-express
- Testing: Jest con supertest
- Linting: ESLint

## Struttura del Progetto
notes-api/
+-- _CONTEXT.md
+-- package.json
+-- .env.example
+-- .gitignore
+-- src/
|   +-- index.js
|   +-- app.js
|   +-- routes/
|   |   +-- notes.js
|   +-- controllers/
|   |   +-- notesController.js
|   +-- services/
|   |   +-- notesService.js
|   +-- middleware/
|   |   +-- errorHandler.js
|   |   +-- validateRequest.js
|   +-- schemas/
|   |   +-- noteSchema.js
|   +-- utils/
|       +-- apiResponse.js
+-- tests/
    +-- setup.js
    +-- notes.test.js

## Modello Dati (Nota) -- storage in-memory
- id: string UUID v4
- title: string (1-200 caratteri)
- content: string (1-10000 caratteri)
- tags: string[] (opzionale, default [])
- createdAt: string ISO 8601
- updatedAt: string ISO 8601

## Endpoint API
GET    /api/notes       Lista tutte le note      200
GET    /api/notes/:id   Dettaglio nota            200/404
POST   /api/notes       Crea nota                 201
PUT    /api/notes/:id   Aggiorna nota             200/404
DELETE /api/notes/:id   Elimina nota              204/404
GET    /api/health      Health check              200
GET    /api/docs        Swagger UI                200

## Formato Risposte
Successo: { "success": true, "data": T }
Con lista: { "success": true, "data": T[], "count": N }
Errore:   { "success": false, "error": { "message": "", "code": "" } }

## Convenzioni
- ESModules: usa SEMPRE import/export, MAI require
- Async: usa SEMPRE async/await, MAI callback o .then()
- camelCase per variabili/funzioni, PascalCase per classi
- kebab-case per nomi file
- Status code appropriati (201 POST, 204 DELETE)

## Vincoli
- NON usare require(). Progetto ESModules.
- NON mettere logica di business nei controller. Usa i service.
- NON ritornare stack trace negli errori.
- NON hardcodare la porta. Usa process.env.PORT default 3000.
- CORS abilitato per lo sviluppo.
- Ogni endpoint deve avere commento JSDoc Swagger.
- package.json DEVE avere "type": "module"
\end{lstlisting}

% ---------------------------------------------------------
\section{Generazione del Progetto}

\begin{praticabox}[title={Pratica --- Inizializzazione e generazione}]
In Copilot Agent Mode:
\begin{lstlisting}[language={}]
Leggi il _CONTEXT.md e implementa l'intero progetto Notes API.

Inizia con:
1. Inizializza npm con package.json (type: module)
2. Installa le dipendenze necessarie
3. Implementa i file nell'ordine:
   schemas -> utils -> services -> middleware ->
   controllers -> routes -> app -> index
4. Aggiungi la configurazione Swagger
5. Crea i test di integrazione
6. Crea .env.example e .gitignore
\end{lstlisting}
\end{praticabox}

% ---------------------------------------------------------
\section{Verifica Strutturale}

\subsection*{Checklist di revisione}

\textbf{\texttt{package.json}:}
\begin{itemize}
  \item Contiene \texttt{"type": "module"}
  \item Script \texttt{dev}, \texttt{start}, \texttt{test}, \texttt{lint} presenti
  \item Dipendenze corrette
\end{itemize}

\textbf{\texttt{src/app.js}:}
\begin{itemize}
  \item Usa \texttt{import}, non \texttt{require}
  \item CORS abilitato
  \item Error handler come ultimo middleware
  \item Route \texttt{/api/notes} e \texttt{/api/docs} montate
\end{itemize}

\textbf{\texttt{src/schemas/noteSchema.js}:}
\begin{itemize}
  \item Schema Zod con validazione title (1--200), content (1--10000)
  \item Campo tags opzionale
\end{itemize}

\textbf{\texttt{src/utils/apiResponse.js}:}
\begin{itemize}
  \item Formato risposta conforme al contesto
  \item Funzioni \texttt{successResponse()} e \texttt{errorResponse()}
\end{itemize}

% ---------------------------------------------------------
\section{Test e Prima Esecuzione}

\begin{praticabox}[title={Pratica --- Avvio del server}]
\begin{lstlisting}[language=bash]
npm run dev
\end{lstlisting}
Output atteso:
\begin{lstlisting}[language={}]
Server running on port 3000
Swagger docs available at http://localhost:3000/api/docs
\end{lstlisting}
\end{praticabox}

\begin{praticabox}[title={Pratica --- Test manuali}]
\textbf{Health check:}
\begin{lstlisting}[language=bash]
curl http://localhost:3000/api/health
\end{lstlisting}
\begin{lstlisting}[language=json]
{ "success": true, "data": { "status": "ok" } }
\end{lstlisting}

\textbf{Crea una nota:}
\begin{lstlisting}[language=bash]
curl -X POST http://localhost:3000/api/notes \
  -H "Content-Type: application/json" \
  -d '{"title":"La mia prima nota",
       "content":"Generata in 0-code!",
       "tags":["test"]}'
\end{lstlisting}

\textbf{Validazione (errore atteso):}
\begin{lstlisting}[language=bash]
curl -X POST http://localhost:3000/api/notes \
  -H "Content-Type: application/json" \
  -d '{"title":""}'
\end{lstlisting}
\begin{lstlisting}[language=json]
{ "success": false,
  "error": { "message": "Validation failed",
             "code": "VALIDATION_ERROR" } }
\end{lstlisting}
\end{praticabox}

\begin{praticabox}[title={Pratica --- Swagger UI}]
Apri nel browser: \url{http://localhost:3000/api/docs}

Dovresti vedere la documentazione interattiva Swagger con tutti gli endpoint.
\end{praticabox}

\begin{praticabox}[title={Pratica --- Test automatici}]
\begin{lstlisting}[language=bash]
npm test
\end{lstlisting}
\end{praticabox}

\begin{checkpointbox}
Server avviato, health check risponde, CRUD funziona, validazione blocca input invalidi, Swagger UI accessibile, test automatici passano.
\end{checkpointbox}

% ---------------------------------------------------------
\section{Comprendere l'Architettura}

\begin{figure}[H]
\centering
\includegraphics[width=0.7\textwidth]{05-http-request-flow}
\caption{Flusso di una richiesta HTTP}
\end{figure}

Ogni layer ha una responsabilità precisa:
\begin{itemize}
  \item \textbf{Router}: definisce URL e metodi HTTP
  \item \textbf{Middleware}: validazione, autenticazione (Cap.~8)
  \item \textbf{Controller}: gestisce request/response HTTP
  \item \textbf{Service}: logica di business (non sa nulla di HTTP)
  \item \textbf{Storage}: persiste i dati
\end{itemize}

\begin{notebox}
Questa separazione in layer è un pattern architetturale classico. Quando nel Capitolo~7 aggiungeremo PostgreSQL, dovremo modificare \textsc{solo} il layer di storage. Tutto il resto resta invariato.
\end{notebox}

% ---------------------------------------------------------
\section{Preparazione per i Prossimi Capitoli}

Questa API è il fondamento del progetto full-stack:

\begin{center}
\begin{tabularx}{\textwidth}{c X}
\toprule
\textbf{Capitolo} & \textbf{Cosa aggiungiamo} \\
\midrule
Cap.~7 & PostgreSQL + Prisma (sostituiamo in-memory storage) \\
Cap.~8 & Autenticazione OAuth~2.0 + JWT \\
Cap.~9 & Frontend React con autenticazione \\
Cap.~10 & Integrazione Full-Stack completa \\
Cap.~11--12 & App Flutter connessa a questa API \\
\bottomrule
\end{tabularx}
\end{center}

\begin{praticabox}[title={Pratica --- Inizializza Git}]
\begin{lstlisting}[language=bash]
git init
git add .
git commit -m "feat: Notes API REST completa"
\end{lstlisting}
\end{praticabox}

% ---------------------------------------------------------
\section*{Riepilogo del Progetto}

\begin{center}
\begin{tabularx}{\textwidth}{l X}
\toprule
\textbf{Aspetto} & \textbf{Dettaglio} \\
\midrule
Stack & Node.js~20 + Express + Zod + Swagger \\
Endpoint & 6 endpoint REST + documentazione \\
Validazione & Schema Zod su tutti gli input \\
Test & Integration test con Jest + supertest \\
Architettura & Router $\rightarrow$ Controller $\rightarrow$ Service $\rightarrow$ Storage \\
Tempo totale & $\sim$45--60 minuti \\
Codice scritto a mano & Solo il \texttt{\_CONTEXT.md} \\
\bottomrule
\end{tabularx}
\end{center}
